\documentclass[a4paper,11pt]{article}

% !TEX program = lualatex

\usepackage{luatexja}
\usepackage{luatexja-fontspec}
\setmainjfont{HaranoAjiMincho}
\setsansjfont{HaranoAjiGothic}

\usepackage{fontspec}
\usepackage{amsmath,amssymb,amsthm}
\usepackage{hyperref}

\title{エルブラン空間における命題の不在と四句分別}
\author{千葉 将臣}
\date{2026-02-19}

\newtheorem{proposition}{命題}
\newtheorem{remark}{注}

\begin{document}
\maketitle

\begin{abstract}
本稿は、エルブラン空間において命題$P$の「不在」を、単なる充足不能（解なし）判定とは異なる仕方で表現するために必要となる除去条件を整理する。
その際、$P$に関して自然に現れる四つの形式$P,\neg P, P\wedge \neg P, \neg(P\vee \neg P)$が、四句分別（catuṣkoṭi）の構成要素として理解できることを指摘する。
\end{abstract}

\section{導入}
エルブラン空間（Herbrand universe/Herbrand base）に基づく推論では、命題の充足可能性や導出可能性が、基礎項の組合せに対する評価として扱われる。
しかし、ある命題$P$が「存在しない」あるいは「表現されていない」ことを示したい場面では、単に「$P$が解なしである（充足不能）」と言うだけでは意図が尽くされないことがある。
本稿の目的は、解なし判定以外で$P$の不在を示すために、エルブラン空間からどの形式を除去すべきかを、最小限の要請として定式化することである。

さらに本稿の動機として、仏教哲学において「空」を導出（あるいは概念的に確定）する際に四句分別が用いられることが、単なる教説上の選択ではなく、ある種の論理的必然性（少なくとも形式的必然性）を持つことを示したい。
すなわち、エルブラン空間における命題の除去操作が四句分別の枠組みを要請するという主張を通じて、四句分別が「否定の徹底」や「不在の言明」を行う上で自然に現れる形式であることを論理学的側面から裏づける。

\section{前提と用語}
以下では、命題$P$に対して、その否定$\neg P$、両立$P\wedge \neg P$、および排中律の否定$\neg(P\vee \neg P)$を考える。
これらは古典論理では互いに単純な関係にあるが、エルブラン空間上での「述語の出現／不出現」や、推論系の表現力の観点では、区別して扱う必要がある。

加えて、「エルブラン空間に（何かを）追加する」という表現は、厳密には次のいずれか（または組合せ）を指し得るため、以後の議論では区別して用いる。
\begin{enumerate}
\item \textbf{シグネチャへの追加}：関数記号・定数記号・述語記号（原子式を作るための語彙）を新たに導入すること。
\item \textbf{エルブラン宇宙への追加}：定数や関数記号の追加により、基礎項（ground term）の集合が拡大すること。
\item \textbf{エルブラン基底への追加}：述語記号の追加や、項の増加に伴い、基礎原子（ground atom）の集合が拡大すること。
\item \textbf{理論（プログラム）への追加}：公理・節・ルールを追加して、許されるエルブラン解釈（モデル）や導出関係を制約すること。
\end{enumerate}
本稿で「$P$を（エルブラン空間から）除去する／追加する」と言う場合、主として(1)--(3)の語彙的・構文的な操作を念頭に置くが、(4)の操作と区別しない限り、命題の真偽は容易に揺れ得る。

\paragraph{「除去」の暫定定義}
本稿では差し当たり、\emph{除去}を「シグネチャ（したがってエルブラン基底）から当該の式形成可能性を奪う語彙的操作」として捉える。
すなわち、言語$\mathcal{L}$からあるパターンの原子・複合式を生成できないように、述語記号（あるいは否定・連言・選言といった形成規則）を制限した言語$\mathcal{L}'$へ移行することを「除去」と呼ぶ。
この意味での除去は、特定の解釈で偽にする（あるいは理論が導出しない）こととは異なり、当該の形式が\emph{そもそも言えない}状態を目標とする。

\paragraph{解なし判定に関する方法論的注意（短絡のリスク）}
本稿は、解なし判定（充足不能性）の利用自体を禁じるものではない。
ただし、解なし判定を得た時点で検証を打ち切り、$P$の「不在」を結論してしまうこと（以下\emph{短絡}と呼ぶ）は、空を外在化してしまう空見へ滑るリスクを伴う。
したがって本稿では、短絡を規範的に禁止するというよりも、短絡がどのように空見の構図へ接続されるかを可視化し、四句分別に対応する四形式の除去を最後まで追う必要性を論理学的に位置づける。

\begin{remark}
なお、短絡の「誘惑」（少ない情報で結論へ飛ぶ圧力）を定量化する一案として、解なし判定と四句分別の検証とで記述長や計算量を比較する枠組みが考えられる。
例えばコルモゴロフ複雑性$K(\cdot)$を用いて、$K(\text{解なし判定}) \ll K(\text{四句分別の検証})$のような不均衡が見込まれるとき、短絡リスクが高いとみなす、といった形式化である。
ただし本稿では、この方向性の具体化は今後の課題として留保する。
\end{remark}

\section{主要主張}
\begin{proposition}[不在の完全除去条件]
エルブラン空間において、解なし判定以外の仕方で命題$P$の不在を示すには、エルブラン空間から
\[
P,\quad \neg P,\quad P\wedge \neg P,\quad \neg(P\vee \neg P)
\]
の四つの形式を除去しなければ、$P$を完全に除去したことにはならない。
\end{proposition}

\begin{proposition}[四句分別としての解釈]
上の四つの形式
\[
P,\quad \neg P,\quad P\wedge \neg P,\quad \neg(P\vee \neg P)
\]
は、四句分別（catuṣkoṭi）を構成する。
\end{proposition}

\section{考察}
命題$P$の「不在」を表現する際、$P$そのものだけでなく、$P$に関する否定・矛盾・両否定の形がエルブラン空間内に残存していると、推論はなお$P$に関する情報を取り出し得る。
したがって、解なし判定に依存しない不在の主張を成立させるためには、四句分別に対応する四形式を一括して除去することが要請される、という見方が得られる。

\subsection{解なし判定による「除去」と空見}
解なし判定（充足不能性）によって、ある理論$T$のもとで命題$P$が成立しえないことを示すことはできる。
しかしこの操作は、前節で暫定定義した意味での「除去」（すなわち、$P$に関する形式が\emph{そもそも言えない}状態を作ること）とは異なる。
というのも、解なし判定は、言語における$P$の形成可能性を残したまま、外部から（公理・制約として）不整合性や不充足を与えることで$P$を「潰す」からである。

この差異は、仏教哲学における「空」の誤った捉え方として語られる「空見」の構図に対応づけて理解できる。
すなわち、対象の側の構造として空を見いだすのではなく、「外から」否定性（不成立性）を付与して空を措定してしまうならば、空は対象の性質というより観察者側の付加物として現れてしまう。
本稿の観点から言えば、解なし判定はまさに、言語内に$P$の形を残しつつ外部制約で不成立を示す操作であり、空を外在化してしまう空見のモデルとして読める。

\paragraph{補足：対象言語の局所的除去とメタ言語の足場}
「何かを言うこと」は、エルブラン空間（宇宙・基底）そのものの元ではなく、メタ言語レベルの営みである。
本稿が扱う「$P$の四句の除去」は、対象言語のシグネチャから特定の命題$P$を生成する能力を奪う局所的な操作（Local Removal）である。
ここで重要なのは、$P$に関する四句を対象言語から完全に除去したとしても、それを「除去した」と論じるメタ言語の足場や、他の語彙を用いた論理空間そのものが破壊されるわけではない点である。
もし仮にエルブラン空間から全シグネチャを消去するような全体消去（Global Removal）を想定すれば、理論を比較検討する論理的基盤そのものが失われるという別の問題が生じる。しかし、特定の命題$P$の不在を論じる上では、この局所的な四句の除去に留めることで、メタ言語の崩壊を招くことなく「完全な不在」を安全かつ形式的に扱うことができるのである。

\begin{remark}
本稿は、上記二命題の論文化（骨格の提示）を目的とし、エルブラン空間・推論規則・否定の扱い（古典／直観主義／パラ一貫など）の差異に応じた厳密化は今後の課題として残す。
\end{remark}

\section{結論}
エルブラン空間における$P$の不在を、充足不能とは別の概念として扱うならば、$P$に付随する四句分別の四形式を同時に除去する必要がある。
この観点は、論理的表現の「不在」を検討する際の最小要件として位置づけられる。

\end{document}